\documentclass[11pt,a4paper]{article}

\usepackage[margin=1in]{geometry}
\usepackage{amsmath,amssymb,amsfonts,bm,mathtools}
\usepackage{booktabs}
\usepackage{hyperref}
\usepackage{enumitem}

\hypersetup{
  colorlinks=true,
  linkcolor=blue,
  urlcolor=blue,
  citecolor=blue
}

\newcommand{\R}{\mathbb{R}}
\newcommand{\Normal}{\mathcal{N}}
\newcommand{\Id}{I}
\newcommand{\diag}{\mathrm{diag}}

\title{\textbf{Current Offline Wind IDE-Advection Model}}
\author{}
\date{\today}

\begin{document}
\maketitle

\section{Purpose}

This note describes the \emph{current} offline model implemented in the codebase after the shift away from alternating optimization and away from Kalman-style initialization during offline training.

The key design choice is:
\begin{quote}
At forecast time origin, the model uses the \textbf{current observed truth} as the current state, and only learns how to map that state to future states.
\end{quote}

Therefore the current offline model is best understood as a \textbf{deterministic forecast model with a residual uncertainty head}, not as a latent-state filter that must infer an unobserved $x_t$ before predicting $x_{t+1}$.

\section{State and Observation}

We have three stations and two wind components per station. The state is
\begin{equation}
Y_t =
\begin{bmatrix}
u_t(s_1) & v_t(s_1) & u_t(s_2) & v_t(s_2) & u_t(s_3) & v_t(s_3)
\end{bmatrix}^{\!\top}
\in \R^6.
\end{equation}

For offline training and validation, the observation at the current forecast origin is treated as the current state:
\begin{equation}
X_t \equiv Y_t.
\end{equation}

In other words, the current model does \emph{not} train by first inferring a separate latent $p(X_t \mid Y_{0:t})$ with learnable $p_0$ and $m_0$. Those parameters remain in the code only for legacy compatibility, and they are not meant to dominate the current offline objective.

\section{NWP-driven Dynamics Outputs}

Given the recent NWP tensor
\begin{equation}
\mathcal{N}_{t-L+1:t} \in \R^{L \times C \times H \times W},
\end{equation}
the advection network produces the following quantities:
\begin{enumerate}[label=\arabic*.]
  \item a paired mean advection vector
  \begin{equation}
  \mu_t \in \R^4,
  \qquad
  \mu_t =
  \begin{bmatrix}
  \mu_{uu,x}(t) & \mu_{uu,y}(t) & \mu_{vv,x}(t) & \mu_{vv,y}(t)
  \end{bmatrix}^{\!\top},
  \end{equation}
  \item a joint covariance
  \begin{equation}
  \Sigma_t \in \R^{4 \times 4},
  \end{equation}
  \item directional base scales
  \begin{equation}
  s_t =
  \begin{bmatrix}
  s_{\parallel}(t) & s_{\perp}(t)
  \end{bmatrix}
  \in \R_{+}^2,
  \end{equation}
  \item component-wise transport gates
  \begin{equation}
  g_t =
  \begin{bmatrix}
  g_u(t) & g_v(t)
  \end{bmatrix}
  \in [0,g_{\max}]^2,
  \end{equation}
  \item an additive state-bias residual
  \begin{equation}
  b_t \in \R^6.
  \end{equation}
\end{enumerate}

The intended interpretation is:
\begin{itemize}
  \item $\mu_t$ gives the preferred transport direction;
  \item $\Sigma_t$ gives uncertainty broadening of that transport;
  \item $s_{\parallel}(t), s_{\perp}(t)$ define the \emph{deterministic base shape} of the kernel;
  \item $g_t$ controls how much the model deviates from persistence;
  \item $b_t$ lets NWP directly correct systematic amplitude and bias errors.
\end{itemize}

\section{Pairwise Transport Kernels}

Let $h_{ab} \in \R^2$ denote the projected spatial displacement from source site $s_b$ to target site $s_a$.

For each target/source component pair $(i,j) \in \{u,v\}^2$, the model constructs a Gaussian-like transport kernel
\begin{equation}
K_t^{(ij)}(a,b)
\propto
\left|D_t^{(ij)}\right|^{-1/2}
\exp\!\left(
-
\bigl(h_{ab} - d_t^{(ij)}\bigr)^\top
\left(D_t^{(ij)}\right)^{-1}
\bigl(h_{ab} - d_t^{(ij)}\bigr)
\right),
\end{equation}
followed by row normalization over source site index $b$.

The drift term $d_t^{(ij)}$ is derived from the corresponding blocks of $\mu_t$, and the covariance term has the form
\begin{equation}
D_t^{(ij)} = B_t^{(ij)} + 2 \Delta t^2 \, \Sigma_{t,\mathrm{pair}}^{(ij)} + \varepsilon \Id_2.
\end{equation}

Here:
\begin{itemize}
  \item $B_t^{(ij)}$ is a directional base covariance obtained from $(s_{\parallel}(t), s_{\perp}(t))$ and the local transport direction;
  \item $\Sigma_{t,\mathrm{pair}}^{(ij)}$ is the relevant pairwise block induced by the joint covariance $\Sigma_t$;
  \item $\varepsilon \Id_2$ is only a numerical stabilizer.
\end{itemize}

This means the kernel width is no longer a hard-coded $\Id_2$ floor. Instead:
\begin{equation}
\text{kernel width} = \text{deterministic base shape} + \text{uncertainty broadening}.
\end{equation}

\section{Gated State Transition}

The normalized pairwise kernels are assembled into a block transport operator
\begin{equation}
\widetilde{A}_t \in \R^{6 \times 6}.
\end{equation}

Let
\begin{equation}
\rho_t = 1 - \Delta t \, \lambda_t,
\end{equation}
where $\lambda_t$ is the learned damping coefficient. The damped transport target is
\begin{equation}
T_t = \rho_t \widetilde{A}_t.
\end{equation}

Define the state-level gate matrix by repeating the component gates over sites:
\begin{equation}
G_t =
\diag\!\bigl(
g_u(t), g_v(t),
g_u(t), g_v(t),
g_u(t), g_v(t)
\bigr).
\end{equation}

The actual transition used by the current model is
\begin{equation}
A_t = \Id_6 + G_t \bigl(T_t - \Id_6 \bigr).
\end{equation}

This is the most important structural change. It implies:
\begin{equation}
g_t = 0
\quad \Longrightarrow \quad
A_t = \Id_6,
\end{equation}
so zero transport now means exact persistence rather than automatic spatial smoothing.

The one-step deterministic forecast is
\begin{equation}
\widehat{Y}_{t+1 \mid t} = A_t Y_t + b_t.
\end{equation}

\section{Residual Uncertainty Model}

Although the offline forecast path is deterministic, the model still keeps scalar process and observation noise parameters:
\begin{equation}
q_{\mathrm{proc}}(t), \qquad r_{\mathrm{obs}}(t).
\end{equation}

They are \emph{not} used to infer the current state during offline training. Instead, they parameterize the predictive residual scale for one-step teacher-forced forecasts:
\begin{equation}
E_t = Y_{t+1} - \widehat{Y}_{t+1 \mid t},
\end{equation}
\begin{equation}
E_t \sim \Normal\!\left(0, \sigma_t^2 \Id_6\right),
\qquad
\sigma_t^2 = q_{\mathrm{proc}}(t)^2 + r_{\mathrm{obs}}(t+1)^2.
\end{equation}

So in the current formulation:
\begin{itemize}
  \item $q_{\mathrm{proc}}$ is a forecast uncertainty term;
  \item $r_{\mathrm{obs}}$ is a residual observation/error floor;
  \item neither is supposed to replace the transport dynamics.
\end{itemize}

\section{Joint Offline Training Objective}

Training is now performed in a \textbf{single joint stage}: the advection network and the IDE parameters $(q_{\mathrm{proc}}, r_{\mathrm{obs}}, \lambda)$ are optimized together.

For a window
\begin{equation}
Y_{t_0:t_0+K} = (Y_{t_0}, Y_{t_0+1}, \dots, Y_{t_0+K}),
\end{equation}
define the teacher-forced one-step forecasts
\begin{equation}
\widehat{Y}_{t+1 \mid t} = A_t Y_t + b_t,
\qquad
t=t_0,\dots,t_0+K-1.
\end{equation}

Define the one-step MSE:
\begin{equation}
\mathcal{L}_{\mathrm{step}}
=
\frac{1}{K}
\sum_{k=0}^{K-1}
\left\|
\widehat{Y}_{t_0+k+1 \mid t_0+k} - Y_{t_0+k+1}
\right\|_2^2.
\end{equation}

Choose a forecast origin $t_\star$ inside the window. Starting from the \emph{true} state $Y_{t_\star}$, the model rolls out open-loop:
\begin{equation}
\widehat{Y}_{t_\star+h+1 \mid t_\star}
=
A_{t_\star+h}\widehat{Y}_{t_\star+h \mid t_\star}+b_{t_\star+h}.
\end{equation}

This gives the rollout loss
\begin{equation}
\mathcal{L}_{\mathrm{roll}}
=
\frac{1}{H}
\sum_{h=0}^{H-1}
\left\|
\widehat{Y}_{t_\star+h+1 \mid t_\star} - Y_{t_\star+h+1}
\right\|_2^2.
\end{equation}

The one-step Gaussian residual NLL is
\begin{equation}
\mathcal{L}_{\mathrm{nll}}
=
\frac{1}{K}
\sum_{k=0}^{K-1}
\frac{1}{2}
\sum_{d=1}^{6}
\left[
\frac{\bigl(Y_{t_0+k+1}^{(d)}-\widehat{Y}_{t_0+k+1 \mid t_0+k}^{(d)}\bigr)^2}{\sigma_{t_0+k}^2}
\;+\;
\log \sigma_{t_0+k}^2
\;+\;
\log(2\pi)
\right].
\end{equation}

Finally, the total joint objective is
\begin{equation}
\mathcal{L}
=
\alpha_{\mathrm{step}}\mathcal{L}_{\mathrm{step}}
\;+\;
\alpha_{\mathrm{roll}}\mathcal{L}_{\mathrm{roll}}
\;+\;
\alpha_{\mathrm{nll}}\mathcal{L}_{\mathrm{nll}}
\;+\;
\alpha_{\mathrm{noise}}\mathcal{R}_{\mathrm{noise}}
\;+\;
\alpha_{\mathrm{smooth}}\mathcal{R}_{\mathrm{smooth}}
\;+\;
\alpha_{\mathrm{gate}}\mathcal{R}_{\mathrm{gate}}
\;+\;
\alpha_{\mathrm{bias}}\mathcal{R}_{\mathrm{bias}},
\end{equation}
where:
\begin{align}
\mathcal{R}_{\mathrm{noise}} &= q_{\mathrm{proc}}^2 + r_{\mathrm{obs}}^2, \\
\mathcal{R}_{\mathrm{smooth}} &= \sum_t \|\theta_{t+1}-\theta_t\|_2^2,
\quad
\theta_t := (\mu_t, \Sigma_t, s_t, g_t, b_t), \\
\mathcal{R}_{\mathrm{gate}} &= \sum_t \|g_t\|_2^2, \\
\mathcal{R}_{\mathrm{bias}} &= \sum_t \|b_t\|_2^2.
\end{align}

This objective explicitly encodes the modeling preference:
\begin{quote}
Stay close to persistence unless transport or NWP residual correction produces a real forecast gain.
\end{quote}

\section{Validation Protocol}

The current validation should match the current training semantics:
\begin{enumerate}[label=\arabic*.]
  \item \textbf{Teacher-forced one-step validation:} use the true current state $Y_t$ to predict $Y_{t+1}$.
  \item \textbf{Open-loop multi-step validation:} use the true state at forecast origin and then recursively roll forward without future truth updates.
\end{enumerate}

Persistence remains the main baseline:
\begin{equation}
\widehat{Y}_{t+h}^{\mathrm{pers}} = Y_t.
\end{equation}

If the proposed model cannot beat persistence under this protocol, then either:
\begin{enumerate}[label=\arabic*.]
  \item the transport correction is still too weak,
  \item the residual bias path is misused,
  \item or the state-transition family is still too restrictive.
\end{enumerate}

\section{What Is No Longer the Main Offline Story}

The following objects still exist in the codebase for compatibility, but they are no longer the intended center of the offline model:
\begin{itemize}
  \item a learned initial prior mean $m_0$,
  \item a learned initial covariance scale $p_0$,
  \item alternating IDE-only versus ML-only optimization rounds.
\end{itemize}

The current offline formulation should instead be read as:
\begin{quote}
\textbf{current truth} $+$ \textbf{gated transport correction} $+$ \textbf{NWP residual bias correction} $+$ \textbf{residual uncertainty calibration}.
\end{quote}

\end{document}
